\section{Limitations}
\label{sec:limitations}

The main experiments use GPT-2 small and a deliberately planted
actor-name cue. This setup makes the cue controllable, but it limits
generalization to larger models and to naturally occurring
correlations. We capped cleaned movie reviews at
\reviewTokenCapPilot{} tokens to limit compute. Before truncation,
\trainOverCapPctPilot{}\% of training reviews and
\testOverCapPctPilot{}\% of test reviews exceeded this cap.
Truncation can remove sentiment evidence, so the relative influence
of the cue may differ on full reviews. The audit recorded
\duplicateReviewPairsPilot{} identical train/test review pairs.
This overlap limits the independence of the affected test examples.

The preregistered trigger-restoration follow-up (Phase~4b), the
Pythia retraining experiment and the finer Pythia swap windows were
not run within the project's time. The partly reliable name-cue
experiment (Phase~L6) stopped after its pilot because no reliability
level produced a qualifying shortcut. Consequently, the results do not
establish how trigger restoration would behave, whether cue strength
explains the rank-dependent pattern, or whether the retraining findings
transfer to Pythia. The read-out fix was tested on the full fine-tuned
GPT-2 models only, its strength was chosen at the edge of the
pre-registered grid, and the layer it uses was identified from earlier
results on the same models. Three seeds make all-seed rules fragile, so
we report per-seed values throughout. The film-versus-movie cue (Phase~5) did not meet the
preregistered learning criterion, including after the prevalence
fallback. Its downstream interventions were therefore not run, and
the name-cue findings remain untested in that setting.
